\chapter{Exoplanets}

\section*{Theory problems}

\probhead{9.1}{}{2010}{Beijing, China}{TH S7}{10}
% source id: 2010_TH_S7   tag: Earth-like exoplanet transit depth
We are interested in finding habitable exoplanets. One way to achieve this is
through the dimming of the star, when the exoplanet transits across the stellar
disk and blocks a fraction of the light. Estimate the maximum luminosity ratio
change for an Earth-like planet orbiting a star similar to the Sun.

\probhead{9.2}{}{2011}{Katowice, Poland}{TH L1}{30}
% source id: 2011_TH_L1   tag: HD209458 transit + Doppler: star radius
A transit of duration 180 minutes was observed for a planet which orbits the
star HD\,209458 with a period of 84 hours. The Doppler shift of absorption lines
arising in the planet's atmosphere was also measured, corresponding to a
difference in radial velocity of $30$\,km\,s$^{-1}$ (with respect to
observer) between the beginning and the end of the transit. Assuming a circular
orbit exactly edge-on to the observer, find the approximate radius and mass of
the star and the radius of the orbit of the planet.

\probhead{9.3}{}{2012}{Rio de Janeiro, Brazil}{TH S11}{}
% source id: 2012_TH_S11   tag: Sun's astrometric wobble from Jupiter
What is the angular amplitude of the oscillatory motion of the Sun, due to the
existence of Jupiter, as measured by an observer located at Barnard's Star? What
is the period of this oscillation?

\probhead{9.4}{Exomoon}{2017}{Phuket, Thailand}{TH T13}{60}
% source id: 2017_TH_T13   tag: exomoon TTV/TDV, Hill sphere, Roche limit
Exomoons are natural satellites of exoplanets. The gravitational influence of
such a moon will affect the position of the planet relative to the planet--moon
barycentre, resulting in Transit Timing Variations ($\sigma_{TTV}$, TTVs) as the
observed transit of the planet occurs earlier or later than the predicted time
of transit for a planet without a moon.

The motion of the planet around the planet--moon barycentre will also induce
Transit Duration Variations ($\sigma_{TDV}$, TDVs) as the observed transit
duration is shorter or longer than the predicted transit duration for a planet
without a moon.

We will consider edge-on circular orbits with the following parameters:
\begin{itemize}\itemsep2pt
\item $M_p$ is the planet mass
\item $M_m$ is the moon mass
\item $P_p$ is the planet--moon barycentre's period around the host star
\item $P_m$ is the moon's period around the planet
\item $a_p$ is the distance of the planet--moon barycentre to the star
\item $a_m$ is the distance of the moon to the planet--moon barycentre
\item $f_m$ is the moon phase, $f_m = 0^\circ$ when the moon is in opposition to
  the star
\item $\tau$ is the mean transit duration of the planet (as if it has no moon)
\end{itemize}

We will only consider the orbit of a prograde moon orbiting in the same plane as
the planet's orbit. Example phases of the moon, as observed by distant
observers, are shown in the figure below.

\ioaafig{2017_TH_T13-f1.pdf}

\begin{description}
\item[(a)] We define $\sigma_{TTV} \equiv t_m - t$ where $t$ is the predicted
  transit time without the moon, and $t_m$ is the observed transit time with the
  moon. Show that
  \[ \sigma_{TTV} = \left[\frac{a_m M_m P_p}{2\pi a_p M_p}\right]\sin(f_m). \]
  A positive value of $\sigma_{TTV}$ indicates that the transit occurs later
  than the predicted time of transit for a planet without a moon. \hfill[10]
\item[(b)] Similarly, we define $\sigma_{TDV} \equiv \tau_m - \tau$ where $\tau$
  is the predicted transit duration without the moon, and $\tau_m$ is the
  observed transit duration with the moon. We can assume that the planet's
  velocity around the star is much bigger than the moon's velocity around the
  planet--moon barycentre, and also that the moon does not change phase during
  the transit. Show that
  \[ \sigma_{TDV} = \tau\left[\frac{P_p M_m a_m}{P_m M_p a_p}\right]\cos(f_m). \]
  A positive value of $\sigma_{TDV}$ indicates that the transit duration is
  longer than the predicted transit duration without a moon. \hfill[13]
\end{description}

An exoplanet is observed transiting a main-sequence solar-type star
($1\,M_\odot$, $1\,R_\odot$, spectral class: G2V). The planet has an edge-on
circular orbit with a period of $3.50$ days. From the observational data, the
planet has a mass of $120\,M_\oplus$ and a radius of $12\,R_\oplus$. The
observed relation between $\sigma_{TTV}^2$ and $\sigma_{TDV}^2$ can be written
as
\[ \sigma_{TDV}^2 = -0.7432\,\sigma_{TTV}^2 + 1.933\times10^{-8}\ \mathrm{days}^2. \]

\begin{description}
\item[(c)] Assume that the moon's mass is much smaller than the planet's mass.
  Find the mean transit duration of the planet ($\tau$) in days. \hfill[6]
\item[(d)] Find the moon's period ($P_m$) in days. \hfill[7]
\item[(e)] Estimate the distance of the moon to the planet--moon barycentre
  ($a_m$) in units of Earth radii. Also find the moon mass ($M_m$) in units of
  Earth mass. \hfill[7]
\item[(f)] The Hill sphere is a region around a planet within which the planet's
  gravity dominates. The radius of the Hill sphere can be written as
  \[ R_h = a_p \sqrt[3]{\frac{M_p}{x M_*}}, \]
  where $M_*$ is the host star mass. Find the value of the constant $x$
  (\textit{Hint:} for a massive host star, the radius of the Hill sphere of the
  system is approximately equal to the distance between the planet and the
  Lagrange point $L_1$ or $L_2$). Hence, find the radius of the Hill sphere of
  this planetary system in units of Earth radii. \hfill[11]
\item[(g)] The Roche limit is the minimum orbital radius at which a satellite
  can orbit without being torn apart by tidal forces. Take the Roche limit as
  \[ R_r = 1.26\,R_p \sqrt[3]{\frac{\rho_p}{\rho_m}}, \]
  where $\rho_p$ and $\rho_m$ are the density of the planet and moon,
  respectively and $R_p$ is the planet's radius. Assuming that the moon is a
  rocky moon with the same density as the Earth, find the Roche limit of the
  system. \hfill[3]
\item[(h)] Does the moon have a stable orbit? (YES / NO) \hfill[3]
\end{description}

\probhead{9.5}{Earth's Transit Zone}{2017}{Phuket, Thailand}{TH T2}{10}
% source id: 2017_TH_T2   tag: Earth-transit zone geometry
Earth's transit zone is an area where extrasolar observers (located far away
from the Solar System) can detect the Earth transiting across the Sun. For
observers on the Earth, this area is the projection of a band around the Earth's
ecliptic onto the celestial plane (light grey area in the left figure). Assume
that the Earth has a circular orbit of 1\,au.

\ioaafig{2017_TH_T2-f1.pdf}

\begin{description}
\item[(a)] Find the angular width of that part of the Earth's transit zone in
  degrees, in which the extrasolar observers can detect Earth's total transit
  (when the whole of the Earth's disk passes in front of the Sun). \hfill[5]
\item[(b)] Find the angular width of that part of the Earth's transit zone in
  degrees, where the extrasolar observers can detect at least Earth's grazing
  transit (when any part of the Earth's disk passes in front of the Sun).
  \hfill[5]
\end{description}

\probhead{9.6}{Life on Other Planets}{2017}{Phuket, Thailand}{TH T7}{20}
% source id: 2017_TH_T7   tag: habitability, MS lifetime, RV detection
One place to search for life is on planets orbiting main sequence stars. A good
starting point is the planets that have an Earth-like temperature range and a
small temperature fluctuation. Assume that for a main sequence star, the
relation between the luminosity $L$ and the mass $M$ is given by
$L \propto M^{3.5}$. You may assume that the total energy $E$ released over the
lifetime of the star is proportional to the mass $M$ of the star. For the Sun,
it will have a main sequence lifetime of about 10 billion years. The stellar
spectral types are given in the table below. Assume that the spectral subclasses
of stars (0--9) are assigned on a scale that is linear in $\log M$.

\begin{center}
\begin{tabular}{lccccccc}
\toprule
Spectral Class & O5V & B0V & A0V & F0V & G0V & K0V & M0V \\
\midrule
Mass ($M_\odot$) & 60 & 17.5 & 2.9 & 1.6 & 1.05 & 0.79 & 0.51 \\
\bottomrule
\end{tabular}
\end{center}

\begin{description}
\item[(a)] If it takes at least $4\times10^{9}$ years for an intelligent life
  form to evolve, what is the spectral type (accurate to the subclass level) of
  the most massive star in the main sequence around which astronomers should
  look for intelligent life? \hfill[6]
\item[(b)] Assume that the target planet has the same emissivity $\epsilon$ and
  albedo $a$ as the Earth. In order to have the same temperature as the Earth,
  express the distance $d$, in au, of the planet to its parent main sequence
  star, of mass $M$. \hfill[6]
\item[(c)] The existence of a planet around a star can be shown by the variation
  in the radial velocity of the star about the star--planet system centre of
  mass. If the smallest Doppler shift in the wavelength detectable by the
  observer is $(\Delta\lambda/\lambda) = 10^{-10}$, calculate the lowest mass
  of such a planet in (b), in units of Earth masses, that can be detected by
  this method, around the main sequence star in (a). \hfill[8]
\end{description}

\probhead{9.7}{Terrestrial Transit}{2021}{Bogot\'a (remote)}{TH Q10}{15}
% source id: 2021_TH_Q10   tag: transit duration geometry
\textit{Note:} Assume perfect circular orbits in both questions below.

\begin{description}
\item[(10.1)] An alien astronomer from a distant planetary system is observing
  the Sun. Suddenly, the brightness of the Sun drops due to the transit of the
  Earth in front of it. What is the maximum duration that this transit may last
  (in hours)? Assume that the planet where the astronomer observes from, does not
  move relative to the Sun. \hfill[5.0\,pt]
\item[(10.2)] Imagine that the transit of a given exoplanet as seen from Earth
  lasts 31 minutes. The host star is a red dwarf, with mass and radius that are
  10\% of the mass and radius of the Sun. What is the minimum orbital period
  this exoplanet may have (in days)? \hfill[10.0\,pt]
\end{description}

\probhead{9.8}{Macondo and Melquiades}{2021}{Bogot\'a (remote)}{TH Q6}{12}
% source id: 2021_TH_Q6   tag: exoplanet RV, orbital radius and period
In 2019, as a part of the NameExoWorlds campaign of the International
Astronomical Union, Colombia was granted an opportunity to select a name for the
star HD\,93083 and its planetary system. HD\,93083 is a $K$-type dwarf star and
has one extrasolar planet orbiting it. Today they are officially known as
Macondo (star) and Melquiades (planet), from the literary ideas of the Colombian
writer Gabriel Garc\'ia M\'arquez.

This star has an effective temperature of $4995$\,K and an apparent visual
magnitude of $8.3$. As per GAIA DR2, the parallax for Macondo is $35.03$
milliarcseconds. You may assume the orbit of Melquiades is perfectly circular.
In the figure you can see the plot of radial velocity of Macondo with respect to
the phase.

\ioaafig{2021_TH_Q6-f1.pdf}

Radial velocity of Macondo (Y-axis in $km\ s^{-1}$) as a function of the phase
(X-axis).

\begin{description}
\item[(6.1)] Find the wavelength (in nm) of peak emission for Macondo in its
  rest frame (i.e., ignoring Doppler shifts). \hfill[2.0\,pt]
\item[(6.2)] Find the distance of this system from the Earth (in parsecs) and
  the absolute visual magnitude ($M_V$) of the star. \hfill[2.0\,pt]
\item[(6.3)] Calculate the mean radial velocity of Macondo (in km\,s$^{-1}$).
  \hfill[2.0\,pt]
\item[(6.4)] Calculate the orbital velocity (in km/s) of Melquiades
  ($v_p$), if mass of the star ($m_s$) is $0.7\,M_\odot$ and the mass of
  exoplanet ($m_p$) is $7\times10^{26}$\,kg. Assume that the orbital plane of
  the system is edge-on with respect to our line-of-sight. \hfill[2.0\,pt]
\item[(6.5)] Find the orbital radius of Melquiades (in au) and its orbital
  period (in days). \hfill[4.0\,pt]
\end{description}

\probhead{9.9}{Circumbinary planet}{2022}{Kutaisi, Georgia}{TH 2}{10}
% source id: 2022_TH_2   tag: circumbinary planet synodic period, illumination
A stellar system consists of two main-sequence stars each of $2M_\odot$ orbiting
with a period of 4 years in a circular orbit. A circumbinary planet orbits the
center of the binary system in the same orbital plane and direction at a fixed
distance of 20\,au.

\begin{description}
\item[(a)] Calculate the amount of time it takes for the planet to reappear at
  the same position relative to the binary system. \hfill[6\,pt]
\item[(b)] At a given instant, what is the maximum fraction of the planet's
  surface area which receives light from at least one of the stars? Neglect any
  atmospheric effects and the sizes of the stars. \hfill[4\,pt]
\end{description}

\section*{Data-analysis problems}

\probhead{9.10}{Stars and Exoplanets}{2021}{Bogot\'a (remote)}{DA DA2}{75}
% source id: 2021_DA_DA2   tag: transit parameters, habitable zone
In this problem, we will explore the connection between the physical properties
of exoplanets and their host stars and will use the observational data to
discover as much as possible about these systems. You may neglect interstellar
extinction.

\medskip
\noindent\textbf{Part 1 (20 points).}

\begin{center}
\begin{tabular}{cccccc}
\toprule
Name of & Name of & $T_{\mathrm{eff}}$ & $g$ & $m_v$ & parallax \\
planet & star & (K) & (m\,s$^{-2}$) & (magnitudes) & (milliarcsec) \\
\midrule
Gorgona & HD\,209458 & 5980 & 347 & 7.63 & 20.67 \\
\bottomrule
\end{tabular}

\smallskip
Table 1: Observational data for exoplanet Gorgona and its parent star HD\,209458
\end{center}

The effective temperature ($T_{\mathrm{eff}}$) and the gravitational
acceleration in the surface of the star ($g$) can be measured from the shape of
the spectrum and its absorption lines.
% sic: "in the surface" as in the original
The visual apparent magnitude ($m_v$) and parallax are measured by doing
photometry and astrometry, respectively.

Additionally, it has been observed that every $3.52$ days the brightness of the
star drops due to the transit of the planet in front of it, as it is
represented in this lightcurve:

\ioaafig{2021_DA_DA2-f1.pdf}
% f1 = Figure 1: Normalized Flux (y-axis) against time (x-axis) for the parent
% star HD 209458

\begin{description}
\item[(1.1)] Use the given information to calculate the following quantities
  for the HD\,209458 system: \hfill[20\,pt]

  \begin{center}
  \begin{tabular}{ccccc}
  \toprule
  Luminosity of & Radius of & Mass of & Mean planet's & Radius of the planet\\
  the star & the star & the star & orbital radius & in Jupiter's radius\\
  $L_\star$ $[L_\odot]$ & $R_\star$ $[R_\odot]$ & $M_\star$ $[M_\odot]$
    & $a$ $[\mathrm{au}]$ & $R_p$ $[R_J]$\\
  \midrule
  \phantom{xxxx} & \phantom{xxxx} & \phantom{xxxx} & \phantom{xxxx}
    & \phantom{xxxx}\\
  \bottomrule
  \end{tabular}
  \end{center}

  \textbf{Note:} Assume that the bolometric correction for all F and G type
  stars is the same.
\end{description}

\medskip
\noindent\textbf{Part 2 (25 points).}

The habitable zone is defined as the region in which a planet may have liquid
water on its surface. This is mainly related to the amount of radiation
received from the host star, which must be within some limits to ensure a
favorable range of planet surface temperatures.

We define the effective flux received by a planet as
$S_{\mathrm{eff}} = \frac{L}{a^2}$, where $L$ is the star luminosity in solar
units, and $a$ is the mean orbital radius in au. The minimum flux in the
habitable zone can be approximated by
\[
S_{\mathrm{min}} = S_{\mathrm{eff}_\odot} + n\cdot T_\star + b\cdot T_\star^2
  + c\cdot T_\star^3 + d\cdot T_\star^4,
\]
where $T_\star = (T_{\mathrm{eff}} - T_{\mathrm{eff}_\odot})$, and
$S_{\mathrm{eff}_\odot}$ is the equivalent flux for the case of the Sun, which
along with the coefficients $n$, $b$, $c$, $d$ is given in the following
table. The maximum flux for habitability, $S_{\mathrm{max}}$, is found with the
same equation but different constants:

\begin{center}
\begin{tabular}{ccc}
\toprule
Constant & $S_{\mathrm{max}}$ & $S_{\mathrm{min}}$\\
\midrule
$S_{\mathrm{eff}_\odot}$ & $1.0512$ & $0.3438$\\
$n$ & $1.3242\times10^{-4}$ & $5.8942\times10^{-5}$\\
$b$ & $1.5418\times10^{-8}$ & $1.6558\times10^{-9}$\\
$c$ & $-7.9895\times10^{-12}$ & $-3.0045\times10^{-12}$\\
$d$ & $-1.8328\times10^{-15}$ & $-5.2983\times10^{-16}$\\
\bottomrule
\end{tabular}
\end{center}

The table below gives real data for 7 different star-planet systems. However
planet names have been changed to honour some natural sanctuaries in Colombia:

\begin{center}
\begin{tabular}{cc cc}
\toprule
\multicolumn{2}{c}{Stellar Parameters} & \multicolumn{2}{c}{Planet Parameters}\\
$T_{\mathrm{eff}}$ [K] & $M_V$ [mag] & Name & $a$ [au]\\
\midrule
6180 & 3.68 & Tayrona   & 0.04\\
5730 & 3.87 & Iguaque   & 0.04\\
5980 & 4.21 & Gorgona   & 0.04\\
5480 & 6.04 & Amacayacu & 0.08\\
5770 & 3.48 & Malpelo   & 0.05\\
6130 & 3.07 & Pisba     & 0.03\\
6140 & 3.85 & Tatam\'a  & 0.06\\
\bottomrule
\end{tabular}
\end{center}
% \ioaafig{2021_DA_DA2-f2.pdf} % f2 = image of the 7-system table transcribed
% just above; kept here (commented) to avoid printing the same table twice

\begin{description}
\item[(2.1)] In the following figure, the vertical axis represents the
  effective temperature of stars, and the horizontal axis represents the
  effective flux received by orbiting planets. The dot marked on the graph
  represents planet Earth, and dashed lines mark the limits of the habitable
  zone. \hfill[15\,pt]

  % NOTE: the original shows here an unlabelled grid (T_eff versus S_eff) with
  % the Earth marked as a dot and the habitable-zone limits as dashed lines;
  % no figure was extracted for it in the crop set (f1-f5).

  Put numerical labels at tickmark position on both axes. Draw on the same
  figure the exact position where Gorgona and Amacayacu would be, as if they
  were also located at 1\,au from their corresponding stars.
\item[(2.2)] Now considering the real orbital radius given in the table for
  each planet, indicate with YES or NO which of them are in the habitable
  zone. Show your quantitative reasoning on the working
  sheets. \hfill[10\,pt]

  \begin{center}
  \begin{tabular}{cc}
  \toprule
  Planet's Name & In habitable zone? YES / NO\\
  \midrule
  Tayrona   & \\
  Iguaque   & \\
  Gorgona   & \\
  Amacayacu & \\
  Malpelo   & \\
  Pisba     & \\
  Tatam\'a  & \\
  \bottomrule
  \end{tabular}
  \end{center}
  % \ioaafig{2021_DA_DA2-f3.pdf} % f3 = image of the blank YES/NO answer table
  % transcribed just above; kept here (commented) to avoid duplication
\end{description}

\medskip
\noindent\textbf{Part 3 (30 points).}

In the last page you find a list of 38 exoplanets, and the goal is to find out
if low-mass exoplanets (LME) and high-mass exoplanets (HME) tend to orbit
around stars with different characteristics.

\begin{description}
\item[(3.1)] To get a robust low-mass subsample one can apply a technique
  called ``iterative sigma-clipping''. The idea is to compute the mean ($\mu$)
  and the standard deviation ($\sigma$) of the masses and to exclude from the
  sample, those planets with masses above $\mu + \sigma$. Then repeat the same
  steps with the remaining subsample two more times. We will say that planets
  in the final subsample are the low-mass ones, and those excluded during the
  iterations, the high-mass ones. Fill the following table with the numbers
  you find in the process: \hfill[10\,pt]

  \begin{center}
  \begin{tabular}{cccccc}
  \toprule
  Low-mass sample & Sample size & $\mu$ & $\sigma$ & $\mu+\sigma$
    & No.\ of planets\\
  & & & & & to exclude\\
  \midrule
  Full / Original & 38 & & & & \\
  subsample after 1st iteration & & & & & \\
  subsample after 2nd iteration & & & & & \\
  subsample after final iteration & & --- & --- & --- & ---\\
  \bottomrule
  \end{tabular}
  \end{center}
  % \ioaafig{2021_DA_DA2-f4.pdf} % f4 = image of the blank sigma-clipping
  % answer table transcribed just above; kept here (commented) to avoid
  % duplication
\item[(3.2)] Make a plot using the X-axis for the serial number of the planets
  in the list $(1, 2, 3, \ldots)$, and the Y-axis for the mass of the
  planets. Mark 3 horizontal lines at the $\mu + \sigma$ thresholds you found
  in the iterations. \hfill[5\,pt]
\item[(3.3)] Let's investigate the possible difference in the effective
  temperatures of host stars in both groups, computing some descriptive
  statistics: \hfill[10\,pt]

  \begin{center}
  \begin{tabular}{cccccc}
  \toprule
  & Min. & 1st Quartile & Median (2nd quartile) & 3rd Quartile & Max\\
  \midrule
  LME & & & & & \\
  HME & & & & & \\
  \bottomrule
  \end{tabular}
  \end{center}
\item[(3.4)] Draw boxplots summarizing the numbers you just computed. Do you
  see a clear difference in the temperatures of the stars orbited by low-mass
  and high-mass planets? Write YES or NO. \hfill[5\,pt]

  % NOTE: the original provides here an empty chart frame labelled
  % "Planetary type vs host star temperature (K)", with HME and LME on the
  % vertical axis and temperature ticks 5200-6400 K on the horizontal axis;
  % no figure was extracted for it in the crop set (f1-f5).
\end{description}

\medskip
\noindent The list of 38 exoplanets (planet mass in Jupiter masses,
$T_{\mathrm{eff}}$ of the host star in K):

\begin{center}
\begin{tabular}{clcc}
\toprule
No. & Name of Planet & Planet Mass $[M_J]$ & $T_{\mathrm{eff}}$ of the Host Star\\
\midrule
1 & KEPLER-37 b & 0.01 & 5520\\
2 & KEPLER-21 b & 0.02 & 6256\\
3 & HD 97658 b  & 0.02 & 5468\\
4 & HD 46375 b  & 0.23 & 5345\\
5 & HD 219134 h & 0.28 & 5209\\
6 & HD 88133 b  & 0.30 & 5582\\
7 & HD33283 b   & 0.33 & 5877\\
8 & HD 149026 b & 0.36 & 6096\\
9 & BD-10 3166 b & 0.46 & 5578\\
10 & HD 75289 b & 0.47 & 6196\\
\bottomrule
\end{tabular}
\end{center}
% TABLE ABRIDGED — full data in crop edition
% (remaining rows: 11 HD 217014 b 0.47 5755; 12 HD 2638 b 0.48 5564;
% 13 WASP-13 b 0.49 6025; 14 WASP-34 b 0.59 5771; 15 HD 209458 b 0.69 5988;
% 16 HAT-P-30 b 0.71 6177; 17 WASP-76 b 0.92 6133; 18 WASP-74 b 0.97 5727;
% 19 HAT-P-6 b 1.06 6442; 20 HD189733 b 1.14 5374; 21 WASP-82 b 1.24 6257;
% 22 KELT-7 b 1.29 6460; 23 HD 149143 b 1.33 6067; 24 KELT-3 b 1.42 6404;
% 25 KELT-2A b 1.49 6164; 26 HD86081 b 1.50 6015; 27 HAT-P-7 b 1.74 6270;
% 28 HD 118203 b 2.14 5847; 29 HAT-P-14 b 2.20 6490; 30 WASP-38 b 2.71 6178;
% 31 HD17156 b 3.20 5985; 32 KELT-6 c 3.71 6176; 33 HD 75732 d 3.86 5548;
% 34 HD 115383 b 4.00 5891; 35 HD 120136 b 5.84 6210; 36 WASP-14 b 7.34 6195;
% 37 HAT-P-2 b 8.74 6439; 38 XO-3 b 11.79 6281)

\ioaafig{2021_DA_DA2-f5.pdf}
% f5 = image of the full 38-planet list (all rows), kept as the complete
% data table

\probhead{9.11}{30 Years of Exoplanets}{2025}{Mumbai, India}{DA D01}{90}
% source id: 2025_DA_D01   tag: radial velocity plus transit, limb darkening
This problem explores some aspects of the two main methods of exoplanet
detection: radial velocity and transit. Throughout this problem we shall
consider a particular system of a single planet (P) in a circular orbit with
radius $a$ around a solar-type star (S). We shall refer to this system as the
``SP system''.

\medskip\noindent\textbf{Stellar parameters}

\begin{description}
\item[(D01.1)] The V-band apparent magnitude of the star S is
  $(7.65\pm0.03)$\,mag, the parallax is $(20.67\pm0.05)$\,milliarcsecond and
  the bolometric correction (BC) is $-0.0650$\,mag. Thus the star has a higher
  bolometric luminosity than its luminosity in the V-band.

  Estimate the mass of the star, $M_{\mathrm{s}}$ (in units of $M_\odot$),
  assuming a mass-luminosity ($M$--$L$) relation of the form $L \propto M^4$.
  Also estimate the uncertainty in $M_{\mathrm{s}}$. You may need
  $\frac{d}{dx}\ln x = \frac{1}{x}$. \hfill[8]
\end{description}

\medskip\noindent\textbf{Radial Velocity method}

The radial velocity method uses the Doppler shift
$\delta\lambda \equiv \lambda_{\mathrm{obs}} - \lambda_0$ between the observed
wavelength $\lambda_{\mathrm{obs}}$ and the rest wavelength $\lambda_0$ of a
known spectral line to detect an exoplanet and determine its characteristics.

The figure below shows the $\delta\lambda$ for the Fe\,I line
($\lambda_0 = 543.45\times10^{-9}$\,m) as a function of time as observed for
the SP system.

\ioaafig{2025_DA_D01-f1.pdf}
% f1 = scatter plot of delta-lambda (in 10^-15 m) against time (0-24 days)

The radial velocity semi-amplitude $K$ is defined as
$K \equiv (v_{\mathrm{r,max}} - v_{\mathrm{r,min}})/2$ where
$v_{\mathrm{r,max}}$ and $v_{\mathrm{r,min}}$ are the minimum and maximum
radial velocities, respectively.
% sic: the original indeed says "the minimum and maximum radial velocities"
% in this order for v_r,max and v_r,min
For a circular planetary orbit the semi-amplitude $K$ can be written as:
\[
K = \left(\frac{2\pi G}{T}\right)^{1/3}
    \frac{M_{\mathrm{p}}\sin i}{(M_{\mathrm{p}} + M_{\mathrm{s}})^{2/3}}
\]
where $T$ is the period, $i$ is the inclination of the planetary orbit (angle
between the normal to the orbital plane of the planet and the line of sight of
the observer), $M_{\mathrm{p}}$ and $M_{\mathrm{s}}$ are the masses of the
planet and the star, respectively.

\begin{description}
\item[(D01.2)] Use the graph given in the Summary Answersheet to answer the
  following.
\item[(D01.2a)] Draw a smooth curve associated with the observed data shown in
  the graph. \hfill[2]
\item[(D01.2b)] Select appropriate points on your drawn curve and use suitable
  methods to determine $T$ and $K$ along with respective uncertainties. All
  data points used for the calculation of $T$ and $K$ must be shown in the
  table in the Summary Answersheet. Use the rest of the table to show your
  intermediate calculations, as needed, with appropriate headers. \hfill[11]
\item[(D01.2c)] Find the minimum mass of the planet $M_{\mathrm{p,min}}$ (in
  $M_\odot$), and its corresponding uncertainty, assuming
  $M_{\mathrm{p}} \ll M_{\mathrm{s}}$. \hfill[5]
\item[(D01.2d)] Using the value of $M_{\mathrm{p,min}}$ estimated in part
  (D01.2c), calculate the minimum value of the semi-major axis of the planet's
  orbit, $a_{\mathrm{min}}$, in au and its uncertainty. \hfill[4]
\end{description}

\medskip\noindent\textbf{Transit method (without limb darkening)}

The schematic diagram of a planet transit (not drawn to scale) is shown
below. Initially, we shall assume the stellar disk to have a uniform average
intensity with some intrinsic noise due to the star itself.

\ioaafig{2025_DA_D01-f2.pdf}
% f2 = schematic: face-on and side-on views of the transit (uniform disk),
% with R_s, bR_s, R_p, inclination i, and the schematic lightcurve showing
% Delta, t_F and t_T

The lightcurve of the normalized intensity, $I$, as a function of time $t$ is
shown in the schematic diagram of the transit above. The average stellar
intensity outside the transit is taken as unity. The maximum decrease in the
intensity is given by $\Delta$ in the normalized light curve. For a uniformly
bright stellar disk, the radius of the planet, $R_{\mathrm{p}}$, is related to
$\Delta$ as
\[
\left(\frac{R_{\mathrm{p}}}{R_{\mathrm{s}}}\right)^2 = \Delta,
\]
where $R_{\mathrm{s}}$ is the radius of the star.

The total duration of transit (when part or all of the planet covers the
stellar disk) is given by $t_{\mathrm{T}}$, while $t_{\mathrm{F}}$ gives the
duration when the planet is fully in front of the stellar disk. The ``impact
parameter'' $b$ is the projected distance between the planet and centre of the
stellar disk at the mid-point of the transit, in units of the stellar radius,
$R_{\mathrm{s}}$.

For a nearly edge-on star-planet orbit, the impact parameter is given by the
formula
\[
b = \left[\frac{(1-\sqrt{\Delta})^2
      - (t_{\mathrm{F}}/t_{\mathrm{T}})^2(1+\sqrt{\Delta})^2}
     {1 - (t_{\mathrm{F}}/t_{\mathrm{T}})^2}\right]^{1/2}
\]

\begin{description}
\item[(D01.3)] For the SP system, the stellar radius is known to be
  $R_{\mathrm{s}} = 1.20\,R_\odot$, and the transit of the planet is indeed
  visible. Using the minimum orbital radius, $a_{\mathrm{min}}$, estimated in
  part (D01.2d), find the minimum value, $i_{\mathrm{min}}$, of the
  inclination angle. \hfill[3]
\end{description}

Assuming a stellar disk of uniform brightness, the transit lightcurve would
look like as shown below.

\ioaafig{2025_DA_D01-f3.pdf}
% f3 = transit lightcurve: normalized intensity (0.980-1.000) against days
% relative to transit midpoint (-0.15 to +0.15)

\begin{description}
\item[(D01.4)] Using the given lightcurve answer the following questions. For
  your reference the above lightcurve is also given in the Summary
  Answersheet.
\item[(D01.4a)] Estimate the values of $t_{\mathrm{T}}$ and $t_{\mathrm{F}}$
  in days by marking appropriate readings on the graph. \hfill[3]
\item[(D01.4b)] Estimate the mean value of $\Delta$, by marking appropriate
  readings on the graph and hence find $R_{\mathrm{p}}$ in units of
  $R_\odot$. \hfill[2]
\item[(D01.4c)] Determine the value of $i$ in degrees assuming the orbital
  radius to be $a_{\mathrm{min}}$. \hfill[2]
\end{description}

\medskip\noindent\textbf{Introducing limb darkening}

So far we have assumed the stellar disk to be uniformly bright. In reality,
the observed brightness of the stellar disk is not uniform due to ``limb
darkening'' --- an optical effect where the central part of the stellar disk
appears brighter than the edge, or the ``limb''.

The limb darkening effect can be measured by the relative intensity
$J(\theta) \equiv \frac{I(\theta)}{I(0)}$, where $\theta$ is the angle between
the normal to the stellar surface at a point and the line joining the observer
to that point, $I(\theta)$ is the observed intensity of the stellar disk at
that point ($I(0)$ being the intensity at the centre of the stellar disk). For
a distant observer, $\theta$ varies from $\theta = 0$ (centre of the disk) to
$\theta \approx 90^\circ$ (edge of the disk).

\begin{description}
\item[(D01.5)] The table below gives measured $J(\theta)$ at a certain
  wavelength for the Sun. We shall assume that the same limb darkening profile
  holds for the star S.

  \begin{center}
  \begin{tabular}{cc}
  \toprule
  $\theta$ & $J(\theta)$\\
  \midrule
  $0^\circ$  & 1.000\\
  $10^\circ$ & 0.994\\
  $15^\circ$ & 0.984\\
  $20^\circ$ & 0.971\\
  $25^\circ$ & 0.950\\
  $30^\circ$ & 0.943\\
  $40^\circ$ & 0.883\\
  $50^\circ$ & 0.794\\
  $60^\circ$ & 0.724\\
  $70^\circ$ & 0.595\\
  $80^\circ$ & 0.475\\
  $90^\circ$ & 0.312\\
  \bottomrule
  \end{tabular}
  \end{center}

  The limb darkening profile can be modelled by a quadratic formula:
  \[
  J(\theta) = 1 - a_1(1-\cos\theta) - a_2(1-\cos\theta)^2,
  \]
  where $a_1$ and $a_2$ are two constants.

  We shall estimate the unknown coefficients $a_1$ and $a_2$ from the given
  data by making a plot with suitable variables.
\item[(D01.5a)] Choose a pair of variables $(x_1, y_1)$ which are suitable
  functions of $\theta$ and $J$, that you want to plot along $x$ and $y$ axes,
  respectively, to determine $a_1$ and $a_2$. Write the expressions for $x_1$
  and $y_1$.

  If you need to define additional variables for additional plots, define them
  as $(x_2, y_2)$, etc. \hfill[2]
\item[(D01.5b)] Tabulate the values necessary for your plots. \hfill[4]
\item[(D01.5c)] Plot the newly defined variables on the given graph paper
  (mark your graph as ``D01.5c''). \hfill[7]
\item[(D01.5d)] Obtain $a_1$ and $a_2$ from the plot. Uncertainties on the
  values are not needed. \hfill[7]
\end{description}

\medskip\noindent\textbf{Transit in the presence of limb darkening}

Now, we consider planetary transits across a limb darkened stellar disk. In
the presence of limb darkening, which we shall model by the quadratic formula
of $J(\theta)$ given above, the average observed intensity of the entire
stellar disk (without any transit), $\langle I\rangle$, is given by:
\[
\langle I\rangle = \left(1 - \frac{a_1}{3} - \frac{a_2}{6}\right) I(0)
\]
Further, the dip in the light caused by the transiting planet now depends not
only on the relative size of the planet and the star,
$\left(\frac{R_{\mathrm{p}}}{R_{\mathrm{s}}}\right)$, but also on the
intensity profile of the stellar disk along the transit chord, which in turn,
depends on the angle of inclination, $i$.

The schematic diagram below (not drawn to scale) shows the configuration. Note
that the brighter part of the star is shown in a darker shade, while the
planet is shown as a black dot.

\ioaafig{2025_DA_D01-f4.pdf}
% f4 = schematic: face-on view (transit chord through point C) and side-on
% view (angle theta_C between line of sight and surface normal) of a transit
% across a limb-darkened disk

Here the relation between
$\left(\frac{R_{\mathrm{p}}}{R_{\mathrm{s}}}\right)$ and the measured $\Delta$
from the light curve is
\[
\Delta = \frac{I(\theta_{\mathrm{C}})}{\langle I\rangle}
  \left(\frac{R_{\mathrm{p}}}{R_{\mathrm{s}}}\right)^2,
\]
where $I(\theta_{\mathrm{C}})$ is the intensity of the stellar disk at the
midpoint of the transit chord (point C in the figure above),
$\theta_{\mathrm{C}}$ being the angle between the line of sight and the normal
to the surface at that point. From the above it is obvious that for a given
star, the same value of $\Delta$ can be produced by many combinations of the
planet size, $R_{\mathrm{p}}$, and the inclination angle $i$.

\begin{description}
\item[(D01.6)] It is possible to uniquely determine both $R_{\mathrm{p}}$ and
  $i$ by using data from transit lightcurves at two wavelengths, say,
  $\lambda_{\mathrm{B}}$ (blue) and $\lambda_{\mathrm{R}}$ (red). The limb
  darkening coefficients for these two wavelengths are given below:

  \begin{center}
  \begin{tabular}{ccc}
  \toprule
  Wavelength & $a_1$ & $a_2$\\
  \midrule
  $\lambda_{\mathrm{B}}$ & 0.82 & 0.05\\
  $\lambda_{\mathrm{R}}$ & 0.24 & 0.20\\
  \bottomrule
  \end{tabular}
  \end{center}
\item[(D01.6a)] Choose the correct statement among the following that
  describes the relation between the maximum depth of the transit ($\Delta$)
  for $\lambda_{\mathrm{B}}$ and the inclination angle ($i$) of the orbit and
  tick ($\checkmark$) it in the Summary Answersheet.

  \begin{itemize}
  \item[(A)] $\Delta$ increases with decreasing $i$.
  \item[(B)] $\Delta$ decreases with decreasing $i$.
  \item[(C)] $\Delta$ is independent of $i$.
  \end{itemize}
  \hfill[2]
\item[(D01.6b)] The maximum depth of the transit ($\Delta$) for the ``SP
  system'' was measured to be 0.0182 and 0.0159 for $\lambda_{\mathrm{B}}$ and
  $\lambda_{\mathrm{R}}$, respectively.

  Draw schematic transit light curves for both $\lambda_{\mathrm{B}}$ and
  $\lambda_{\mathrm{R}}$ on the given grid and label the curves by ``B'' and
  ``R'' respectively. Assume that the total transit duration is same for both
  wavelengths. The curves need not be to scale, but should represent the
  shapes of the light curves correctly. \hfill[4]
\item[(D01.7)] We shall use a graphical method to find the values of
  $R_{\mathrm{p}}$ and $i$ for the SP system using measurements of $\Delta$ at
  $\lambda_{\mathrm{B}}$ and $\lambda_{\mathrm{R}}$.
\item[(D01.7a)] Write an appropriate expression connecting the relevant
  variables that are to be plotted. (Hint: You may consider $i$ or $b$, and
  $R_{\mathrm{p}}$ among the relevant variables.) \hfill[6]
\item[(D01.7b)] Tabulate the appropriate quantities that are to be
  plotted. \hfill[5]
\item[(D01.7c)] Draw a suitable graph and mark it as ``D01.7c''. \hfill[7]
\item[(D01.7d)] Estimate the values of $R_{\mathrm{p}}$ (in $R_\odot$) and $i$
  (in degrees) from the graph. \hfill[4]
\item[(D01.8)] Based on the results obtained in this problem, indicate whether
  the planet P is ``ROCKY'' or ``GASEOUS'' by ticking ($\checkmark$) the
  appropriate box in the Summary Answersheet. \hfill[2]
\end{description}

\section*{Observation --- planetarium}

\probhead{9.12}{TRAPPIST-1}{2023}{Katowice, Poland}{OBS-PLAN 3}{35}
% source id: 2023_OBS-PLAN_3   tag: planetarium exoplanet transits, resonances
Aliens have found out that Earth's astronomers discovered planets in the
TRAPPIST-1 system by observing numerous transits. They have used their flying
saucer (similar to the one you were in for the observation round) to take you
to the 5th planet (designated~f) of TRAPPIST-1, and have asked you to show
them the methods Earthlings use to uncover the parameters of the system. A
clock displaying time in Earth hours will be visible. The whole presentation
lasts 520\,h (1\,s represents 1\,h). (Projection time 10 minutes)

Based on your observations (you can use the space on the last sheet for
observing notes),
\begin{description}
\item[(a)] determine the following quantities for the planet you are on (use
  Earth hours for the times): \hfill(7 points)
  \begin{description}
  \item[i.] length of the sidereal day [h]
  \item[ii.] orbital period [h]
  \item[iii.] length of the `solar' day [h]
  \item[iv.] circular orbit --- YES / NO
  \item[v.] obliquity (axial tilt)
  \end{description}
\item[(b)] and the following quantities for each planet b, c, d and e:
  synodic period [h] and maximum elongation [$^\circ$]. \hfill(16 points)
\item[(c)] calculate the orbital period in hours and the semi-major axis in
  tau (where 1 tau $=$ ``TRAPPIST-1f astronomical unit'' $=$ the semi-major
  axis of the orbit of TRAPPIST-1f) of each planet b, c, d, e.
  \hfill(8 points)
\item[(d)] The term `gravitational resonance' is used to describe the
  phenomenon when ratio % sic: source omits "the"
  of the orbital periods of two planets in a system
  is close to the ratio of two integers. Find which pair(s) of planets
  correspond to each of the listed resonances if any: 3:2, 8:5, 5:3, 8:3,
  4:1, 6:1. \hfill(4 points)
\end{description}

\section*{Solutions to Chapter 9}

% Solutions are not transcribed. The official solutions for every theory and
% data-analysis problem are reproduced in full in the crop edition,
% ../IOAA_Problems_by_Topic_2007-2025.pdf, under the same problem number.

\solhead{9.1}{}
% source id: 2010_TH_S7
The flux change is proportional to the ratio of their surface areas, i.e.,
\[ F_e / F_{sun} = (R_e / R_{sun})^2 \]
\hfill(5 points)
\[ (R_e / R_{sun})^2 = 8.4\times10^{-5} \approx 10^{-4} \]
\hfill(5 points)

Obviously this difference is extremely small.

\solhead{9.2}{}
% source id: 2011_TH_L1
Calculation of the angular distance on the orbit which the planet travels
during the transit. \hfill(6 points)\\
Since the orbit is circular,
\[ \sin\frac{\alpha}{2} = \frac{\pi t}{T} \]
where $t$ and $T$ are, respectively, the duration of the transit and the
orbital period of the planet. We can leave the angle as $\sin\frac{\alpha}{2}$
as this will be useful later; the numerical value of this
expressiob % sic: source misprint for "expression"
is about $0.112$.

Calculation of the velocity of the planet. \hfill(10 points)\\
The velocity $v$ at the beginning and end of the transit and the measured
difference in velocities $\Delta v$ form an isosceles triangle, with the angle
between the equal sides $= 2\alpha$. Thus
\[ 2v\cdot\sin\frac{\alpha}{2} = \Delta v \]
from which we get a velocity equal to about $134$~km/h.
% sic: source says "km/h"; the value is in km/s
\hfill(2 points)

Calculation of the radius of the orbit. \hfill(2 points)\\
Since we know the velocity on a circular orbit, this is trivial:
\[ 2\pi r = vT \Rightarrow r = \frac{vT}{2\pi} \]
Substituting we get about $6.4\cdot10^{6}$~km. \hfill(2 points)

Calculation of the radius of the star. \hfill(2 points)\\
With the velocity on the orbit and the transit duration, we get the radius of
the star:
\[ 2R = tv \Rightarrow R = \frac{tv}{2} = 7.5\cdot10^{5}\ \mathrm{km}. \]
The star is thus similar to the Sun in size. \hfill(2 points)

Calculation of the mass of the % sic: sentence truncated in source
star:

Given the orbital speed and radius, do this in the usual way:
\[ \frac{v^2}{r} = \frac{GM}{r^2} \]
\hfill(2 points)

thus
\[ M = \frac{v^2 r}{G} \]
The mass is therefore about $1.8\cdot10^{30}$~kg, so again similar to the
mass of the Sun. \hfill(2 points)

\solhead{9.3}{}
% source id: 2012_TH_S11
The oscillation occurs due to the fact that the center of mass (CM) of the
Solar System does not lie
% sic: source has "lieexactly"
exactly at the center of the Sun, but in another position around which the
center of the Sun revolves.

% note: the source uses the astronomical symbol for Jupiter as subscript;
% rendered here as "Jup".
Given $a_{\mathrm{Jup}} = 5.204$~AU, the distance from CM to the center of the Sun
is:
\[ x_{CM} = \frac{\frac{1}{1047}}{1+\frac{1}{1047}}\,5.204 = 0.004966\ AU, \]
\hfill(3 pt: position of CM)

which is % sic: source has "Whichis"
the amplitude of the movement.

If the distance to Barnard's Star is $d_{bar} = 1.834$~pc, we can write, for
the angular amplitude:
\[ \alpha \approx \frac{x_{cm}}{d_{bar}} = 0.002708'' \]
\hfill(3 pt: angular amplitude)

As for the period, it is the period of Jupiter's orbit:
\[ \frac{T_{\mathrm{Jup}}^2}{a_{\mathrm{Jup}}^3} = \frac{T_\oplus^2}{a_\oplus^3}
   = 1\,\frac{year^2}{AU^3}
   \quad\therefore\quad T_{jup} = 11.87\ years \]
\hfill(4 pt: period of oscillation)

Obs: the difference of Jupiter's mass in the Kepler's Third Law represents a
difference in less than 0.01 year at the period of Jupiter.

\solhead{9.4}{Exomoon}
% source id: 2017_TH_T13
\begin{description}
\item[(a)] From centre of mass for planet--moon system, the distance of the
moon to the planet--moon barycenter can be written as,
\[ a_m = \frac{M_p}{M_m}\,a_{pb} \]
\hfill(2.0)

where $a_{pb}$ is distance of the planet to the planet--moon barycentre.

For the observers on the Earth, projection distance of the planet to the
planet--moon barycenter is
\[ a_{proj} = a_{pb}\sin(f_m) \]
\hfill(3.0)

From angular velocity, TTV signal can be calculated from
\begin{align*}
\sigma_{TTV} &\equiv t_m - t\\
\sigma_{TTV} &= \frac{\theta_{trans}}{\omega_p}
\end{align*}
\hfill(1.0)
\[ \sigma_{TTV} = \left.\left(\frac{a_{proj}}{a_p}\right)\middle/
   \left(\frac{2\pi}{P_p}\right)\right. \]
\hfill(2.0)
\[ \sigma_{TTV} = \frac{a_{proj}P_p}{a_p 2\pi} \]
\[ \sigma_{TTV} = \left[\frac{a_m M_m P_p}{2\pi a_p M_p}\right]\sin(f_m) \]
\hfill(2.0)

\item[(b)] The velocity of planet--moon barycenter around the star is
\[ v_p = \frac{2\pi}{P_p}\,a_p \]
\hfill(1.0)

The velocity of planet around the planet--moon barycentre is
\[ v_{pb} = \frac{2\pi}{P_m}\,a_{pb} \]
\hfill(1.0)

Therefore, the transverse velocity of planet around the planet--moon
barycentre is
\begin{align*}
v_{trans} &= -v_{pb}\cos(f_m)\\
v_{trans} &= -\frac{2\pi M_m}{P_m M_p}\,a_m\cos(f_m)
\end{align*}
\hfill(3.0)

Average transit duration can be written as
\[ \tau = \frac{D}{v_p} \]
\hfill(2.0)

where $D$ is the distance of the planet has to cross in order to complete the
transit. The transit duration with exomoon can be written as
\[ \tau_m = \frac{D}{v_p + v_{trans}} \]
\hfill(2.0)

Therefore, the TDV signal is
\[ \sigma_{TDV} \equiv \tau_m - \tau
   = \frac{\tau v_p}{v_p + v_{trans}} - \tau \]
\hfill(1.0)
\[ \sigma_{TDV} = \tau\left[\frac{-v_{trans}}{v_p + v_{trans}}\right] \]
In case $v_p \gg v_{trans}$
\[ \sigma_{TDV} = -\tau\left[\frac{v_{trans}}{v_p}\right] \]
\hfill(2.0)
\[ \sigma_{TDV} = \tau\left[\frac{P_p M_m a_m}{P_m M_p a_p}\right]\cos(f_m) \]
\hfill(1.0)

\item[(c)] From Kepler's third's law, assume that the moon mass is much
smaller than the planet mass, the distance of the planet--moon barycentre to
the star is
\[ P^2 = \frac{4\pi^2}{G(M_* + M_p)}\,a_p^3 \]
\hfill(1.0)
\[ (3.5\times86400)^2
   = \frac{4\pi^2}{6.67\times10^{-11}\times(1.99\times10^{30}
     + 120\times5.98\times10^{24})}\,a_p^3 \]
\[ a_p = 6.75\times10^{9}\ \mathrm{m} \]

The mean transit duration of the planet can be calculated from the transit
duration of exoplanet without exomoon.
\ioaafig{2017_TH_T13-s1.pdf}
\[ \tau = P\,\frac{\alpha}{2\pi} \]
\hfill(1.0)
\[ \tau = \frac{P}{2\pi}\times 2\sin^{-1}\!\left(\frac{R_* + R_p}{a_p}\right) \]
\hfill(2.0)
\[ \tau = \frac{3.5}{\pi}\,\sin^{-1}\!\left(
   \frac{6.96\times10^{8} + 12\times6.38\times10^{6}}{6.75\times10^{9}}\right)
   \ \mathrm{days} \]
\[ \tau = 0.128\ \mathrm{days} \]
\hfill(2.0)

\item[(d)] The TTV signal is 90 degrees out of phase with the TDV signal in
edge-on circular orbit system. Therefore, the relation between the TTV and TDV
signals is,
\[ \sigma_{TDV}^2 = -\left(\frac{2\pi\tau}{P_m}\right)^2\sigma_{TTV}^2
   + \tau^2\left(\frac{a_m M_m P_p}{a_p M_p P_m}\right)^2 \]
\hfill(3.0)

The observed relation between $\sigma_{TTV}^2$ and $\sigma_{TDV}^2$ can be
written as
\[ \sigma_{TDV}^2 = -0.7432\,\sigma_{TTV}^2 + 1.933\times10^{-8}\ \mathrm{days}^2 \]

Therefore
\[ -\left(\frac{2\pi\tau}{P_m}\right)^2 = -0.7432 \]
\hfill(2.0)
\[ P_m = 0.933\ \mathrm{days} \]
\hfill(2.0)

\item[(e)] From Kepler third law, assuming that moon is small and can be
neglected, the distance of the moon to the planet--moon barycentre is
\[ P_m^2 = \frac{4\pi^2}{GM_p}\,a_m^3 \]
\hfill(1.0)
\[ (0.933\times86400)^2
   = \frac{4\pi^2}{6.67\times10^{-11}\times(120\times5.98\times10^{24})}\,a_m^3 \]
\[ a_m = 1.99\times10^{8}\ \mathrm{m} = 31.2 R_\oplus \]
\hfill(2.0)

From the observed relation between $\sigma_{TTV}^2$ and $\sigma_{TDV}^2$, the
moon mass is
\[ \tau^2\left(\frac{a_m M_m P_p}{a_p M_p P_m}\right)^2
   = 1.933\times10^{-8}\ \mathrm{days}^2 \]
\hfill(2.0)
\[ (0.128)^2\left(\frac{(1.99\times10^{8})\times M_m\times(3.5)}
   {(6.75\times10^{9})\times(120)\times(0.933)}\right)^2
   = 1.933\times10^{-8}\ \mathrm{days}^2 \]
\[ M_m = 1.18 M_\oplus \]
\hfill(2.0)

\item[(f)] For planet--moon system, the angular velocity of barycenter of
planet--moon system is
\begin{align*}
F_{centrifugal-bary} &= F_{star-bary}\\
(M_p + M_m)\,\omega_p^2 a_p &= \frac{GM_*(M_p + M_m)}{a_p^2}\\
\omega_p^2 &= \frac{GM_*}{a_p^3}
\end{align*}
\hfill(2.0)

For exomoon at Lagrange point $L_1$ or $L_2$,
\[ F_{centrifugal-moon} = F_{star-moon} \pm F_{planet-moon} \]
\hfill(1.0)
\[ M_m\omega_p^2(a_p \pm a_m)
   = \frac{GM_*M_m}{(a_p \pm a_m)^2} \pm \frac{GM_pM_m}{(a_{pb}+a_m)^2} \]
\hfill(3.0)
\[ \frac{M_*(a_p \pm a_m)}{a_p^3}
   = \frac{M_*}{(a_p \pm a_m)^2} \pm \frac{M_p}{(a_{pb}+a_m)^2} \]

Since $M_p \gg M_m$ and $a_{pb} \ll a_m$
\[ \frac{M_*(a_p \pm a_m)}{a_p^3}
   = \frac{M_*}{(a_p \pm a_m)^2} \pm \frac{M_p}{a_m^2} \]
\hfill(1.0)
\begin{align*}
M_*(a_p \pm a_m)^3 a_m^2 &= M_*a_p^3 a_m^2 \pm M_p a_p^3 (a_p \pm a_m)^2\\
M_*(\pm 3a_p^2 a_m^3 + 3a_p a_m^4 \pm a_m^5)
  &= \pm M_p a_p^3 (a_p \pm a_m)^2
\end{align*}

For $a_p \gg a_m$, $M_*(3a_p^2 a_m^3) \approx M_p a_p^5$ \hfill(1.0)
\begin{align*}
a_m^3 &= \frac{M_p}{3M_*}\,a_p^3\\
a_m &= a_p\sqrt[3]{\frac{M_p}{3M_*}}
\end{align*}
\[ x = 3 \]
\hfill(1.0)

The radius of the Hill sphere of this planetary system is
\[ R_h = a_p\sqrt[3]{\frac{M_p}{3M_*}} \]
\[ R_h = 6.75\times10^{9}\times
   \sqrt[3]{\frac{120\times5.98\times10^{24}}{3\times1.99\times10^{30}}}
   \ \mathrm{m} \]
\[ R_h = 3.33\times10^{8}\ \mathrm{m} = 52.2 R_\oplus \]
\hfill(2.0)

\item[(g)] From Roche limit equation
\[ R_r = 1.26\,R_p\sqrt[3]{\frac{\rho_p}{\rho_m}} \]
\[ R_r = 1.26\times12\times\sqrt[3]{\frac{120/(12)^3}{1}}\,R_\oplus \]
\[ R_r = 6.21 R_\oplus \]
\hfill(3.0)

\item[(h)] \textbf{Stable orbit} because the distance between the moon and
planet--moon barycentre is between Roche limit and Hill sphere radius
($R_r < a_m < R_H$). \hfill(3.0)

\textit{Mark only if students get the answers for e) \& g) \& h). Students
obtain full mark if answer ``YES'' for $R_r < a_m < R_H$, and vice versa.}
\end{description}

\solhead{9.5}{Earth's Transit Zone}
% source id: 2017_TH_T2
\begin{description}
\item[(a)] For Earth's transit, the whole Earth's disk should pass in front of
the Sun.
\ioaafig{2017_TH_T2-s1.pdf}

From $\triangle STA \sim \triangle ETD$: \hfill(1.0)
\begin{align*}
\frac{X}{R_\oplus} &= \frac{a+X}{R_\odot}\\
X &= \frac{aR_\oplus}{R_\odot - R_\oplus}
\end{align*}
\hfill(1.0)

The half angular size of the Earth's Transit Zone with transit can be written
as,
\[ \theta_T = \arcsin\!\left(\frac{R_\oplus}{X}\right)
            = \arcsin\!\left(\frac{R_\odot}{a+X}\right)
            = \arcsin\!\left(\frac{R_\odot - R_\oplus}{a}\right) \]
\hfill(1.0)

\textit{Solution with}
$\theta_T = \arctan\!\left(\dfrac{R_\odot - R_\oplus}{a}\right)
 \approx \dfrac{R_\odot - R_\oplus}{a}$
\textit{is acceptable with full mark.}

The angular size of the Earth's Transit Zone with transit is
\[ 2\theta_T = 2\arcsin\!\left(\frac{R_\odot - R_\oplus}{a}\right) \]
\hfill(1.0)
\[ 2\theta_T = 0.527^\circ \]
\hfill(1.0)

\item[(b)] For Earth's grazing transit, any part of the Earth's disk should
pass in front of the Sun.
\ioaafig{2017_TH_T2-s2.pdf}

From $\triangle SGB \sim \triangle EGC$: \hfill(1.0)
\begin{align*}
\frac{Y}{R_\oplus} &= \frac{a-Y}{R_\odot}\\
Y &= \frac{aR_\oplus}{R_\odot + R_\oplus}
\end{align*}
\hfill(1.0)

The angular size of the Earth's Transit Zone with transit can be written as,
\[ \theta_G = \arcsin\!\left(\frac{R_\oplus}{Y}\right)
            = \arcsin\!\left(\frac{R_\odot}{a-Y}\right)
            = \arcsin\!\left(\frac{R_\odot + R_\oplus}{a}\right) \]
\hfill(1.0)

\textit{Solution with}
$\theta_G = \arctan\!\left(\dfrac{R_\odot + R_\oplus}{a}\right)
 \approx \dfrac{R_\odot + R_\oplus}{a}$
\textit{is acceptable with full mark.}

The angular size of the Earth's Transit Zone with grazing transit is
\[ 2\theta_G = 2\arcsin\!\left(\frac{R_\odot + R_\oplus}{a}\right) \]
\hfill(1.0)
\[ 2\theta_G = 0.537^\circ \]
\hfill(1.0)
\end{description}

\solhead{9.6}{Life on Other Planets}
% source id: 2017_TH_T7
\begin{description}
\item[(a)] Since $E \propto M \rightarrow L\tau \propto M$ and
$L \propto M^{3.5}$,
\[ \tau = \tau_\odot\left(\frac{M_\odot}{M}\right)^{2.5} \]
\hfill(2.0)

The maximum mass of the star should be
\[ M = M_\odot\left(\frac{\tau_\odot}{\tau}\right)^{1/2.5}
     = M_\odot\left(\frac{10^{10}}{4\times10^{9}}\right)^{1/2.5} \]
\hfill(0.5)
\[ M = 1.44 M_\odot \]
\hfill(0.5)

Use of spectral class--$\log M$ linear relation
\[ \frac{10}{(\log 1.6 - \log 1.05)}\times(\log 1.6 - \log 1.44) = 2.5 \]
\hfill(2.0)

which corresponds to the star of type \textbf{F2V} to \textbf{F3V}. \hfill(1.0)

\textit{0.5 mark for the answer ``between F0V and G0V''.}

\item[(b)] The received power by a planet of a radius $r$ and a distance $d$
from the star is
\[ P_{input} = (1-a)\frac{\pi r^2 L}{4\pi d^2} = (1-a)\frac{r^2 L}{4d^2}, \]
where $a$ is the planet's albedo. \hfill(1.0)

The released power of blackbody radiation is
\[ P_{output} = 4\pi r^2\epsilon\sigma T^4, \]
where $\epsilon$ is emissivity. \hfill(1.0)

Since we assume that the planet emits heat as blackbody radiation,
\[ (1-a)\frac{r^2 L}{4d^2} = 4\pi r^2\epsilon\sigma T^4. \]
\hfill(1.0)

The temperatures of the planet and the Earth are similar,
\[ T^4 = \frac{(1-a)L}{16\pi\epsilon\sigma d^2}
       = \frac{(1-a)L_\odot}{16\pi\epsilon\sigma d_\odot^2}. \]
\[ \frac{L}{d^2} = \frac{L_\odot}{d_\odot^2} \]
\hfill(1.0)
\[ \left(\frac{d}{d_\odot}\right)^2 = \frac{L}{L_\odot}
   = \left(\frac{M}{M_\odot}\right)^{3.5} \]
\hfill(1.0)

The distance of the planet to the star is
$d = \left(\frac{M}{M_\odot}\right)^{3.5/2}$~au. \hfill(1.0)

\item[(c)] Angular speed around the center of mass can be derived by Kepler's
law.
\[ T^2 = \frac{4\pi^2 d^3}{G(M+m)} \]
\hfill(1.0)
\[ \omega = \sqrt{\frac{G(M+m)}{d^3}} \]
\hfill(1.0)

According to the centre of mass of the planet and the star,
\begin{align*}
Md_s &= md\\
d_s &= \frac{m}{M}\,d
\end{align*}
\hfill(1.0)

and
\[ v = \omega d_s \]
\hfill(1.0)
\[ v = \omega d\,\frac{m}{M} = \frac{m}{M}\sqrt{\frac{G(M+m)}{d}}
     \approx m\sqrt{\frac{G}{Md}} \]

From
\[ \frac{\Delta\lambda}{\lambda} = \frac{v}{c} = \frac{m}{c}\sqrt{\frac{G}{Md}} \]
\hfill(1.0)

Use of previously derived expression. The distance of the planet to the star
is
\[ d = \left(\frac{M}{M_\odot}\right)^{3.5/2}\ \mathrm{au}. \]
\hfill(1.0)
\begin{align*}
m &= \frac{\Delta\lambda}{\lambda}\,c\sqrt{\frac{Md}{G}}
   = \frac{\Delta\lambda}{\lambda}\,c
     \sqrt{\frac{M\left(\frac{M}{M_\odot}\right)^{3.5/2}\times 1\ \mathrm{au}}{G}}\\
  &= \frac{\Delta\lambda}{\lambda}\,c
     \sqrt{\frac{M\,(1.44)^{3.5/2}\times 1\ \mathrm{au}}{G}}\\
  &= (299792458\ \mathrm{m/s})\times10^{-10}
     \sqrt{\frac{1.44\times1.99\times10^{30}\ \mathrm{kg}\,
       (1.44)^{3.5/2}\times1.50\times10^{11}\ \mathrm{m}}
       {6.67\times10^{-11}\ \mathrm{m^3kg^{-1}s^{-2}}}}
\end{align*}
\hfill(1.0)
\[ m = 3.31\times10^{24}\ \mathrm{kg} = 0.554 M_\oplus \]
\hfill(1.0)
\end{description}

\solhead{9.7}{Terrestrial Transit}
% source id: 2021_TH_Q10
\begin{description}
\item[(10.1)] The maximum duration occurs if Earth passes exactly along the
diameter of the Sun. Now consider the following figure, where Sun rays are
depicted as traveling to the right towards the hypothetical distant observer
(reason why they can be considered parallel):
\ioaafig[1.0]{2021_TH_Q10-s1.pdf}

The arc along the circumference associated to the transit, $2\alpha$, can be
easily found looking at the shaded triangle:
\[ \alpha = \sin^{-1}\!\left(\frac{R_\odot}{R_{orb}}\right) \]
\hfill(2 points)

being $R_{orb}$ the orbital radius of the Earth. Now the time of the transit
can be found by considering the angular velocity of the Earth, for instance by
means of the following proportionality relations:
\[ \frac{t_{tr}}{T_{orb}} = \frac{2\alpha}{2\pi}
   = \frac{\sin^{-1}\!\left(\frac{R_\odot}{R_{orb}}\right)}{\pi} \]
\hfill(2 points)

Substituting with the known values for these quantities we get, i.e.,
\[ t_{tr} \sim 12.97\ \mathrm{h} = 12\mathrm{h}\,58\mathrm{m} \]
\hfill(1 point)

\item[(10.2)] First of all it must be said that the minimum orbital period is
obtained if we assume that the transit occurs along the diameter of the star.
In other cases, the planet would be crossing a shorter path in front of the
star during the same time, meaning that it would have a smaller angular
velocity and therefore a longer period.

That said, it means that we can resort to the same expression of 10.1, yet
this time we need 2 additional elements:
\begin{itemize}
\item Using the small angle approximation:
\[ \sin(\alpha) \sim \alpha \]
\hfill(1 point)
\item Invoking the full expression for the orbital period:
\[ T_{orb} = \frac{2\pi}{\sqrt{GM_*}}\,R_{orb}^{3/2} \]
\hfill(1 point)
\end{itemize}

Combining these results, we get:
\[ \frac{t_{tr}}{T_{orb}}
   = \frac{t_{tr}}{\frac{2\pi}{\sqrt{GM_*}}R_{orb}^{3/2}}
   = \frac{\frac{R_*}{R_{orb}}}{\pi} \]
\hfill(2 points)

And solving for $R_{orb}$:
\[ R_{orb} = \frac{GM_*\,t_{tr}^2}{2^2 R_*^2} \]
\hfill(2 points)

Finally putting this last result back into the expression for the orbital
period:
\[ T = \frac{\pi G M_*}{4R^3}\,t_{tr}^3 \]
\hfill(2 points)

At this point, just a numerical evaluation is needed, though a more elegant
solution can be achieved by means of scaling relations by noting that this
very same expression must be true in the case of the Earth--Sun system as seen
from far away. Having that 31 minutes is the 4\% of 12.94h:
\[ T = 365.25\cdot\frac{0.1}{0.1^3}\cdot 0.04^3
     \sim 2.3\ \textit{d\'ias} % sic: Spanish "days" left in the source
     \sim 2\mathrm{d}\,7\mathrm{h} \]
\hfill(2 points)

These values were inspired by the planetary system Trappist 1.
\end{description}

\solhead{9.8}{Macondo and Melquiades}
% source id: 2021_TH_Q6
\begin{description}
\item[(6.1)] From Wien's law:
\[ \lambda_{max} = \frac{2.898\times10^{-3}\ \mathrm{m\cdot K}}{T_{eff}\,(\mathrm{K})}
   = \frac{2.898\times10^{-3}\ \mathrm{m\cdot K}}{4995\ \mathrm{K}} \]
\[ \lambda_{max} = 5.8\times10^{-7}\ \mathrm{m} = 580.2\ \mathrm{nm} \]
\hfill(2 points)

\item[(6.2)] With the parallax doable % sic: source wording
to find the distance from Earth:
\[ d = \frac{1}{pllx\,(\mathrm{arcsec})} = \frac{1}{0.035\ \mathrm{arcsec}}
     = 28.55\ \mathrm{pc} \]
\hfill(1 point)
\begin{align*}
M_V &= m_v - 5\log(d\,[\mathrm{pc}]) + 5\\
M_V &= 8.3 - 5\log(28.57) + 5 = 6.02
  % sic: source uses 28.55 pc above but 28.57 here
\end{align*}
\hfill(1 point)

\item[(6.3)] \hfill(2 points)\\
\ioaafig{2021_TH_Q6-s1.pdf}
Each division on Y has 0.002~km\,s$^{-1}$, so the maximum radial velocity is
$43.662$~km\,s$^{-1}$ and the minimus % sic: source misprint for "minimum"
is $43.626$~km\,s$^{-1}$. Then, the mean radial velocity of Macondo is
$v_r = 43.644$~km\,s$^{-1}$.

\item[(6.4)] The tangential velocity of Macondo, from the plot, is its
variation from the mean system velocity
\[ v_s = 43.662\ \mathrm{km\,s^{-1}} - 43.644\ \mathrm{km\,s^{-1}}
       = 0.018\ \mathrm{km\,s^{-1}} \]
The masses of the star and the planet are known so it is only needed to find
orbital velocity of Melquiades. But first it is important to convert the mass
of the star to kg.
\[ 0.7M_\odot \times \frac{1.989\times10^{30}\ \mathrm{kg}}{1M_\odot}
   = 1.4\times10^{30}\ \mathrm{kg} \]
\hfill(1 point)
\[ v_p = \frac{m_s}{m_p}\times v_s
       = \frac{1.4\times10^{30}\ \mathrm{kg}}{7\times10^{26}\ \mathrm{kg}}
         \times 0.02\ \mathrm{km\cdot s^{-1}}
       = 40\ \mathrm{km\cdot s^{-1}} \]
\hfill(1 point)

\item[(6.5)] As the motion of the planet is circular, the orbital period is:
\[ T = \frac{2\pi}{v_p}\times a \]
being $a$ the distance to the central star.

With the 3$^{\mathrm{rd}}$ Kepler's law, the orbital period is
$T^2 = \frac{4\pi^2}{Gm_s}\times a^3$.

Then, combining both expressions of $T$ is possible to find $a$
\begin{align*}
\left(\frac{2\pi}{v_p}\times a\right)^{2} &= \frac{4\pi^2}{Gm_s}\times a^3\\
\left(\frac{2\pi}{v_p}\right)^{2}\frac{Gm_s}{4\pi^2} &= a
\end{align*}
\[ a = 7.25\times10^{10}\ \mathrm{m} = 0.48\ \mathrm{au} \]
\hfill(2 points)

With the previous solution is able to find the orbital period:
\[ T = \frac{2\pi}{v_p}\times a
     = \frac{2\pi\times(7.21\times10^{10}\ \mathrm{m})}{36000\ \mathrm{m\cdot s^{-1}}}
     % sic: source gives a = 7.25e10 m above but uses 7.21e10 m here,
     % and v_p = 40 km/s above but 36000 m/s here
     = 1.27\times10^{7}\ \mathrm{s} = 147\ days \]
\hfill(2 points)
\end{description}

\solhead{9.9}{Circumbinary planet}
% source id: 2022_TH_2
Equating gravitational force to centrifugal force, we get:
\begin{align*}
\frac{4\pi^2}{T_b^2}\,a_b &= \frac{G\,2M_\odot}{4a_b^2}\\
\frac{T_b^2}{a_b^3} &= 2\,\frac{4\pi^2}{GM_\odot}\\
\Longrightarrow a_b &= 2\ \mathrm{au}.
\end{align*}
\hfill(2.0 pt)

The period of planet can be calculated from the 3rd Kepler's law of motion
\begin{align*}
\frac{T_p^2}{a_p^3} &= \frac{4\pi^2}{GM_\Sigma} = \frac{4\pi^2}{4GM_\odot}\\
\Longrightarrow T_p &= 44.7\ \mathrm{years}
\end{align*}
\hfill(2.0 pt)

Now,
\begin{align*}
\frac{1}{T_s} &= \frac{1}{T_b} - \frac{1}{T_p}\\
\therefore\ T_s &= \frac{T_b T_p}{T_p - T_b} = 4.4\ \mathrm{years}
\end{align*}
\hfill(2.0 pt)

One star always covers exactly the half of the planet's surface. In the best
case, the stars are located at the maximal angular distance from each other.
So the fraction of illuminated surface is simply given by
\[ n = \frac{1}{2} + \frac{\theta}{2\pi}
     = \frac{1}{2} + \frac{2\arctan\!\left(\frac{a_b}{a_p}\right)}{2\pi}
     = 0.53. \]
\hfill(4.0 pt)

\solhead{9.10}{Stars and Exoplanets}
% source id: 2021_DA_DA2
\noindent\textbf{Part 1.} \hfill(20 points)

\begin{itemize}
\item With the parallax we get the distance:
$D = \frac{1}{parallax} = 48.38$~pc, \hfill(1 point)\\
then we get the absolute magnitude using the distance modulus
$M_V = m_v - 5\log_{10}(D) + 5 = 4.2067$, \hfill(1 point)\\
and finally we go for the luminosity as BC is assumed to be the same for all
F/G stars, $(M_V - M_{V_\odot}) = (M_{bol} - M_{bol_\odot})$ \hfill(2 points)\\
by comparing with the magnitude of the Sun:
$\frac{L}{L_\odot} = 10^{-0.4(M_V - M_{V_\odot})} = 1.759$. \hfill(1 point)
\item We find the stellar radius from the Stephan--Boltzmann
% sic: source spelling "Stephan-Boltzmann"
law, which in solar units leads to
$\frac{T_{eff\,*}^4}{T_{eff\,\odot}^4}
 = \frac{R_\odot^2}{R_*^2}\frac{L_*}{L_\odot}$ \hfill(2 points)\\
i.e.,
$R_*[R_\odot] = \frac{T_{eff\,\odot}^2\sqrt{L_*\,[L_\odot]}}{T_{eff\,*}^2}
 = 1.238$. \hfill(2 points)
\item Now to find the stellar mass consider the surface gravitational
acceleration, which for Gauss' law and spherical symmetry considerations is:
$g = G\frac{M_*}{R_*^2}$, \hfill(1 point)\\
then one solves $M_*$ for and puts it in the right units,
$M_* = 1.944\ [M_\odot]$ \hfill(3 points)
\item Planet's orbital radius comes from
$a = \sqrt[3]{\frac{GMT^2}{4\pi^2}} = 0.0565$~A.U. \hfill(3 points)
\item Finally, checking the dimming due to the transit this is of
approximately 2.6\%, which is give
% sic: source wording "which is give by"
by the ratio between the areas of the planet disk and the stellar disk, i.e.,
$\left(\frac{R_p}{R_*}\right)^2 = 0.026$, \hfill(2 points)\\
leading to $R_p = 1.988\ [R_J]$ \hfill(2 points)
\end{itemize}

\begin{center}
\begin{tabular}{ccccc}
\toprule
Luminosity of & Radius of & Mass of & Mean planet's & Radius of the planet\\
the star & the star & the star & orbital radius & in Jupiter's radius\\
$L_*[L_\odot]$ & $R_*[R_\odot]$ & $M_*[M_\odot]$ & $a$[au] & $R_p[R_J]$\\
\midrule
1.759 & 1.24 & 1.94 & 0.0565 & 1.988\\
\bottomrule
\end{tabular}
\end{center}

Tolerance: of 5\% on the final answers is acceptable, numbers approximated to
the second digit in the partial calculations are OK.

\ioaafig{2021_DA_DA2-s1.pdf}

\medskip
\noindent\textbf{Part 2.} \hfill(25 points)

\noindent\textbf{(2.1)} \hfill(15 points)\\
One first needs to calibrate the axes. The shown circle representing the Earth
gives as a vertical mark in 5778~$^\circ$K,
% sic: source writes "5778 °K"
and a horizontal one at $S_{eff} = 1$. Then one need to retrieve additional
data from one or more of the additional observations:

\begin{itemize}
\item Having the temperature mark for the Sun, one can estimate $S_{min}$,
adding a second mark on the horizontal axis, so the scale of this one can be
calibrated. Then one can proceed backwards, starting from a given value of
$S_{min}$ or $S_{max}$, find the corresponding temperature, so the vertical
axis can be calibrated as well.
\item Alternatively, one can notice that the limiting $S_{max}$ curve crosses
the X axis in $S_{eff} = 1$, and look numerically for the corresponding
temperature of the bottom line, thus calibrating the vertical axis. Then
starting from a given temperature, one calibrates de
% sic: source wording
horizontal one pivoting on $S_{min}$ / $S_{max}$ calculations as explained
before.

There are several ways, some more precise than others. The student should be
as careful as possible (we give small tolerance to the error in the
calibrations of the axes). A tricky part is that the X-axis increases towards
the left, but this is something rigorous students will discover, as there is
no self-consistency otherwise. Intuition may also help if one knows that
Earth is closer to the high-flux end of the Sun's habitable zone.
\item The luminosity of Gorgona's host star was found in Part A. After
finding the luminosity of Amacayacu's one the final plot looks like this:
\end{itemize}

For placing Gorgona in the right position \hfill(2 pt)

\ioaafig{2021_DA_DA2-s2.pdf}

For calculating flux and placing Amacayacu in the right position \hfill(4 pt)

Tolerance for the position of the planets: $\pm0.025$ in $S_{eff}$; $\pm25$
[K] in $T_{eff}$.

\medskip
\noindent\textbf{(2.2)} \hfill(10 points)\\
NONE of the exoplanets lie in habitable zones. They orbit their stars too
close, so the effective flux is many times larger than $S_{max}$. The
long-but-straightforward answer is to perform all the calculations. The
elegant, clever one, is to note that Amacayacu orbits the faintest/coolest
star, and it is the one farthest away, so its effective flux is the absolute
minimum of the whole set. This effective flux was calculated as ${\sim}0.3$
in A), assuming a distance of 1~A.U. With the real distance, 0.08~A.U, it
grows by a factor $\frac{1}{0.08^2}\sim156$, i.e., it has
$S_{eff}\sim46$, but the habitable zone for this range of temperatures can
not exceed $S_{eff}\sim1.1$, which completely closes the case.

\begin{center}
\begin{tabular}{lr}
\toprule
Name & $S_{eff}$\\
\midrule
Tayrona & 1785.994090\\
Iguaque & 1499.270574\\
Gorgona & 1096.175314\\
Amacayacu & 50.794890\\
Malpelo & 1374.231792\\
Pisba & 5568.747040\\
Tatama & 678.730709\\
\bottomrule
\end{tabular}
\end{center}

\medskip
\noindent\textbf{Part 3.}

\noindent\textbf{(3.1)} (3 points per each row; 1 point extra for the only
number in the last row)

\begin{center}
\begin{tabular}{cccccc}
\toprule
Low-mass sample & Sample size & $\mu$ & $\sigma$ & $\mu+\sigma$
  & No.\ of planets to exclude\\
\midrule
Full / Original & 38 & 1.99 & 2.55 & 4.54 & 4\\
1st subsample & 34 & 1.23 & 1.11 & 2.34 & 5\\
2nd subsample & 29 & 0.84 & 0.61 & 1.45 & 5\\
Final subsample & 24 & -- & -- & -- & --\\
\bottomrule
\end{tabular}
\end{center}

\noindent\textbf{(3.2)}
\ioaafig{2021_DA_DA2-s3.pdf}

\noindent\textbf{(3.3)} (Tolerance 5\%)

\begin{center}
\begin{tabular}{cccccc}
\toprule
$T_{eff}$ & Min. & 1st Quartile & Median & 3rd Quartile & Max\\
\midrule
LME & 5209 & 5574.5 & 5932.5 & 6181.75 & 6460\\
HME & 5548 & 5992.5 & 6177.0 & 6255.0 & 6490\\
\bottomrule
\end{tabular}
\end{center}

\noindent\textbf{(3.4)}
\ioaafig{2021_DA_DA2-s4.pdf}

Yes, the boxplots indicate that HMEs tend to get formed around hotter stars.
\hfill(1 point)

\solhead{9.11}{30 Years of Exoplanets}
% source id: 2025_DA_D01
\begin{description}
\item[(D01.1)] \hfill(8 points)\\
Given\\
Parallax: $p = (20.67\pm0.05)$~milliarcsecond\\
V-band apparent magnitude: $m_{\mathrm{V}} = (7.65\pm0.03)$~mag\\
Let:\\
V-band absolute magnitude: $M_{\mathrm{V}}$\\
Bolometric Magnitude: $M_{\mathrm{bol}}$\\
Bolometric Magnitude of the Sun: $M_{\mathrm{bol},\odot} = 4.74$~mag (given)
\begin{align*}
M_{\mathrm{V}} &= m_{\mathrm{V}} + 5(\log p + 1) = 4.227\ \mathrm{mag}\\
M_{\mathrm{bol}} &= M_{\mathrm{V}} + BC\\
  &= m_{\mathrm{V}} + 5(\log p + 1) + \mathrm{BC} = 4.162\ \mathrm{mag}
\end{align*}
Also
\[ M_{\mathrm{bol}} = M_{\mathrm{bol},\odot} - 2.5\log\frac{L_{\mathrm{bol}}}{L_\odot} \]
\[ \therefore\quad L_{\mathrm{bol}}
   = 10^{\frac{M_{\mathrm{bol},\odot}-M_{\mathrm{bol}}}{2.5}}\,L_\odot
   = 6.521\times10^{26}\ \mathrm{W} \]
Then
\[ M_{\mathrm{s}} = \left(\frac{L_{\mathrm{bol}}}{L_\odot}\right)^{1/4} M_\odot
   = 1.14\ \mathrm{M_\odot} \]
Uncertainty estimate:
\begin{align*}
\Delta M_{\mathrm{bol}} &= \sqrt{(\Delta m_{\mathrm{V}})^2
  + \left(\frac{5}{\ln 10}\,\frac{\Delta p}{p}\right)^2}\\
  &= 0.03\ \mathrm{mag}
\end{align*}
Now,
\[ M_{\mathrm{s}} = 10^{\frac{M_{\mathrm{bol},\odot}-M_{\mathrm{bol}}}{10}} \]
\begin{align*}
\therefore\quad \Delta M_{\mathrm{s}}
  &= M_{\mathrm{s}}\,\frac{\Delta M_{\mathrm{bol}}}{10}\,\ln 10\\
  &= 0.01\ \mathrm{M_\odot}
\end{align*}
\[ \boxed{M_{\mathrm{s}} = (1.14\pm0.01)\ \mathrm{M_\odot}} \]

\item[(D01.2a)] \hfill(2 points)\\
\ioaafig{2025_DA_D01-s1.pdf}

\item[(D01.2b)] \hfill(11 points)\\
$T$ Determination:\\
Least count on time axis $= 0.2$~d\\
There are 13 zero crossings (ZC) and 6 periods.\\
Difference between ZC-1 and ZC-13 $= (21.6-0.6) = 21.0$~d
\[ T = \frac{21}{6} = 3.5\ \mathrm{d} \]
Uncertainty in $T$:
\[ \Delta T = \frac{1}{6}\sqrt{(0.2^2+0.2^2)} = 0.05\ \mathrm{d} \]
\[ \boxed{T = (3.50\pm0.05)\ \mathrm{d}} \]

$K$ Determination:
\begin{center}
\begin{tabular}{cccc}
\toprule
Serial No. & Values of Extrema & $2\delta\lambda$ & $K$\\
of Extrema & $(\times10^{-15}\mathrm{m})$ & m & m\,s$^{-1}$\\
\midrule
1 & $-165$ & & \\
2 & 160 & 325 & \\
3 & $-160$ & & \\
4 & 165 & 325 & \\
5 & $-160$ & & \\
6 & 155 & 315 & \\
7 & $-165$ & & \\
8 & 155 & 320 & \\
9 & $-160$ & & \\
10 & 160 & 320 & \\
11 & $-165$ & & \\
\bottomrule
\end{tabular}
\end{center}
% note: the table is cut by the page break in the source PDF after serial
% no. 11; the quoted mean 2(delta lambda) = 321.67e-15 m implies a 12th
% extremum (2 delta lambda = 325e-15 m) on the truncated row.

Least count on $\delta\lambda$ axis $= 5\times10^{-15}$~m\\
Differences in measured extrema ($\delta\lambda$) of alternate pairs are
written in table.
\begin{align*}
2(\delta\lambda)_{\mathrm{Mean}} &= 321.67\times10^{-15}\ \mathrm{m}\\
(\delta\lambda)_{\mathrm{Mean}} &= 160.83\times10^{-15}\ \mathrm{m}\\
K &= \frac{\delta\lambda}{\lambda}\times c\\
  &= \frac{160.83\times10^{-15}\times2.998\times10^{8}}{543.45\times10^{-9}}\\
  &= 88.73\ \mathrm{m\,s^{-1}}
\end{align*}
Uncertainty in $K$:
\begin{align*}
\sigma_{\delta\lambda} &= \sqrt{\frac{1}{N-1}
  \Sigma\left(\delta\lambda-(\delta\lambda)_{\mathrm{Mean}}\right)^2}\\
  &= 2.041\times10^{-15}\ \mathrm{m}\\
\sigma_K &= \frac{\sigma_{\delta\lambda}}{\lambda}\times c\\
  &= \frac{2.041\times10^{-15}\times2.998\times10^{8}}{543.45\times10^{-9}}
   = 1.12\ \mathrm{m\,s^{-1}}
\end{align*}
\[ \boxed{K = (89\pm1)\ \mathrm{m\,s^{-1}}} \]

\item[(D01.2c)] \hfill(5 points)\\
The minimum mass of the planet
$M_{\mathrm{p,min}} \equiv M_{\mathrm{p}}\sin i$ comes from equation of $K$.
\[ K = \left(\frac{2\pi G}{T}\right)^{1/3}
   \frac{M_{\mathrm{p}}\sin i}{(M_{\mathrm{p}}+M_{\mathrm{s}})^{2/3}} \]
We assume $M_{\mathrm{p}} + M_{\mathrm{s}} \approx M_{\mathrm{s}}$.

For $M_{\mathrm{s}} = 1.14\,\mathrm{M_\odot}$, $T = 3.5$~d:
\[ M_{\mathrm{p,min}} \approx K M_{\mathrm{s}}^{2/3}
   \left(\frac{T}{2\pi G}\right)^{1/3}
   = 6.927\times10^{-4}\ \mathrm{M_\odot} \]
Uncertainty in $M_{\mathrm{p,min}}$:
\begin{align*}
\Delta M_{\mathrm{p,min}} &= M_{\mathrm{p,min}}
  \sqrt{\left(\frac{\Delta K}{K}\right)^2
  + \left(\frac{2}{3}\frac{\Delta M_{\mathrm{s}}}{M_{\mathrm{s}}}\right)^2
  + \left(\frac{1}{3}\frac{\Delta T}{T}\right)^2}\\
  &= 0.091\times10^{-4}\ \mathrm{M_\odot}
\end{align*}
\[ \boxed{M_{\mathrm{p,min}} = (6.93\pm0.09)\times10^{-4}\ \mathrm{M_\odot}} \]

\item[(D01.2d)] \hfill(4 points)
\begin{align*}
a_{\min} &= \left(\frac{T^2G(M_{\mathrm{s}}+M_{\mathrm{p}}\sin i)}{4\pi^2}\right)^{1/3}
  = \left(\frac{T^2G(M_{\mathrm{s}}+M_{\mathrm{p,min}})}{4\pi^2}\right)^{1/3}\\
  &= 0.0471\ \mathrm{au}
\end{align*}
Uncertainty in $a_{\min}$:
\begin{align*}
\Delta a_{\min} &= a_{\min}\sqrt{\left(\frac{2}{3}\frac{\Delta T}{T}\right)^2
  + \left(\frac{1}{3}\right)^2\frac{(\Delta M_{\mathrm{s}})^2
    + (\Delta M_{\mathrm{p,min}})^2}{(M_{\mathrm{s}}+M_{\mathrm{p,min}})^2}}\\
  &= 0.0005\ \mathrm{au}
\end{align*}
\[ \boxed{a_{\min} = (0.0471\pm0.0005)\ \mathrm{au}} \]

\item[(D01.3)] \hfill(3 points)
\[ bR_{\mathrm{s}} = a\cos i \quad\Longrightarrow\quad
   i = \cos^{-1}\frac{bR_{\mathrm{s}}}{a} \]
The minimum value of $i$ corresponds to the maximum value of $b(=1)$ and the
minimum value of $a$. Therefore,
\[ i_{\min} = \cos^{-1}\!\left(\frac{b_{\max}R_{\mathrm{s}}}{a_{\min}}\right)
   = \cos^{-1}\!\left(\frac{R_{\mathrm{s}}}{a_{\min}}\right)
   = \cos^{-1}\!\left(\frac{1.20\ \mathrm{R_\odot}}{0.0471\ \mathrm{au}}\right) \]
\[ \boxed{i_{\min} = 83.20^\circ} \]

\item[(D01.4a)] \hfill(3 points)\\
\ioaafig{2025_DA_D01-s2.pdf}
In the figure of the transit curve, least count on time axis $= 0.006$~d, and
on the intensity axis $= 4\times10^{-4}$.\\
The measurements on the left and right edges of the transit for
$t_{\mathrm{T}}$ are $-0.072$~d and $0.072$~d.\\
The measurements on the left and right edges of the transit for
$t_{\mathrm{F}}$ are $-0.048$~d and $0.048$~d.\\
Therefore,
\[ \boxed{t_{\mathrm{T}} = 0.144\ \mathrm{d}} \qquad
   \boxed{t_{\mathrm{F}} = 0.096\ \mathrm{d}} \]

\item[(D01.4b)] \hfill(2 points)\\
The value of the dip is $(0.9860+0.9848)/2 = 0.9854$
\[ \Delta\ \text{measured from the graph is } 1-0.9854 = 0.0146 \]
\[ \therefore\ \frac{R_{\mathrm{p}}}{R_{\mathrm{s}}} = \sqrt{\Delta} = 0.121 \]
\[ \Longrightarrow\ R_{\mathrm{p}} = 0.121 R_{\mathrm{s}} \]
\[ \boxed{R_{\mathrm{p}} = 0.145\ \mathrm{R_\odot}} \]

\item[(D01.4c)] \hfill(2 points)
\[ b = \left[\frac{(1-\sqrt{\Delta})^2
   - (t_{\mathrm{F}}/t_{\mathrm{T}})^2(1+\sqrt{\Delta})^2}
   {1-(t_{\mathrm{F}}/t_{\mathrm{T}})^2}\right]^{1/2} \]
Plugging in values of $\Delta$, $t_{\mathrm{T}}$, and $t_{\mathrm{F}}$, we get
$b = 0.622$
\[ bR_{\mathrm{s}} = a\cos i \]
Putting $a = a_{\min}$,
\begin{align*}
i &= \cos^{-1}\!\left(\frac{bR_{\mathrm{s}}}{a_{\min}}\right)\\
  &= \cos^{-1}\!\left(\frac{0.622\times1.2\ \mathrm{R_\odot}}{0.0471\ \mathrm{au}}\right)
\end{align*}
\[ \boxed{i = 85.77^\circ} \]

\item[(D01.5a)] \hfill(2 points)\\
Define
\[ \boxed{x_1 = 1-\cos\theta} \qquad\text{and}\qquad
   \boxed{y_1 = \frac{J-1}{1-\cos\theta} \equiv \frac{J-1}{x_1}} \]
Then, $y_1 = -a_1 - a_2 x_1$.\\
In this method $a_1$ and $a_2$ can be determined from the intercept and
slope, respectively, of a best fit straight line to a plot of $y_1$ vs $x_1$.

\item[(D01.5b)] \hfill(4 points)
\begin{center}
\begin{tabular}{cccc}
\toprule
$\theta$ & $J(\theta)$ & $x_1$ & $y_1$\\
\midrule
$0^\circ$ & 1.000 & 0.000 & $\times$\\
$10^\circ$ & 0.994 & 0.015 & $-0.395$\\
$15^\circ$ & 0.984 & 0.034 & $-0.470$\\
$20^\circ$ & 0.971 & 0.060 & $-0.481$\\
$20^\circ$ & 0.970 & 0.060 & $-0.497$\\
$25^\circ$ & 0.950 & 0.094 & $-0.534$\\
$30^\circ$ & 0.943 & 0.134 & $-0.425$\\
$40^\circ$ & 0.883 & 0.234 & $-0.500$\\
$40^\circ$ & 0.890 & 0.234 & $-0.470$\\
$50^\circ$ & 0.794 & 0.357 & $-0.577$\\
$60^\circ$ & 0.724 & 0.500 & $-0.552$\\
$70^\circ$ & 0.595 & 0.658 & $-0.616$\\
$80^\circ$ & 0.475 & 0.826 & $-0.635$\\
$90^\circ$ & 0.312 & 1.000 & $-0.688$\\
\bottomrule
\end{tabular}
\end{center}

\item[(D01.5c)] \hfill(7 points)\\
Plot of $y_1$ vs $x_1$ with straight line fit:
\ioaafig[1.0]{2025_DA_D01-s3.pdf}

\item[(D01.5d)] \hfill(7 points)\\
Slope of best fit line $= -0.232$; Intercept $= -0.456$
\[ \boxed{\Longrightarrow\ a_1 = 0.46, \quad a_2 = 0.23} \]

\item[(D01.6a)] \hfill(2 points)\\
As $i$ decreases, the transit chord moves closer to the limb and the
contribution of the midpoint of the transit to the dip in the average
brightness reduces. Therefore, $\Delta$ decreases with decreasing $i$.\\
Correct option: \fbox{B}

\item[(D01.6b)] \hfill(4 points)\\
Since the limb darkening effect is stronger for $\lambda_{\mathrm{B}}$, it
causes \textit{lesser} reduction in the total light at the ingress and egress
and more reduction at the transit centre, than for $\lambda_{\mathrm{R}}$.
Therefore, the curves must intersect.
\ioaafig{2025_DA_D01-s4.pdf}

\item[(D01.7a)] \hfill(6 points)
\begin{align*}
\Delta &= \frac{I(\theta_{\mathrm{C}})}{\langle I\rangle}
  \left(\frac{R_{\mathrm{p}}}{R_{\mathrm{s}}}\right)^2\\
\Longrightarrow\ \left(\frac{R_{\mathrm{p}}}{R_{\mathrm{s}}}\right)^2
  &= \frac{\Delta\langle I\rangle}{I(\theta_{\mathrm{C}})}
\end{align*}
Here, $I(\theta_{\mathrm{C}})$ is a function of $i$. Further,
$\langle I\rangle = \left(1-\frac{a_1}{3}-\frac{a_2}{6}\right)I(0)$.
\ioaafig{2025_DA_D01-s5.pdf}

From the figure,
\[ \sin\theta_{\mathrm{C}} = \frac{bR_{\mathrm{s}}}{R_{\mathrm{s}}} = b
   = \frac{a\cos i}{R_{\mathrm{s}}} \]
\[ \Longrightarrow\ \cos\theta_{\mathrm{C}} = \sqrt{1-b^2}
   = \sqrt{1-\left(\frac{a\cos i}{R_{\mathrm{s}}}\right)^2} \]
Therefore,
\begin{align*}
I(\theta_{\mathrm{C}}) &= J(\theta_{\mathrm{C}})I(0)\\
  &= \left[1 - a_1(1-\cos\theta_{\mathrm{C}})
    - a_2(1-\cos\theta_{\mathrm{C}})^2\right]I(0)\\
  &= \left[1 - a_1\left(1-\sqrt{1-b^2}\right)
    - a_2\left(1-\sqrt{1-b^2}\right)^2\right]I(0)
\end{align*}
Thus,
\[ \left(\frac{R_{\mathrm{p}}}{R_{\mathrm{s}}}\right)
   = \left[\frac{\Delta\langle I\rangle}{I(\theta_{\mathrm{C}})}\right]^{1/2} \]
\[ \left(\frac{R_{\mathrm{p}}}{R_{\mathrm{s}}}\right)
   = \left[\frac{\Delta\left(1-\frac{a_1}{3}-\frac{a_2}{6}\right)}
   {1 - a_1\left(1-\sqrt{1-b^2}\right)
    - a_2\left(1-\sqrt{1-b^2}\right)^2}\right]^{1/2} \]
\[ \left(\frac{R_{\mathrm{p}}}{R_{\mathrm{s}}}\right)
   = \left[\frac{\Delta\left(1-\frac{a_1}{3}-\frac{a_2}{6}\right)}
   {1 - a_1\left(1-\sqrt{1-\left(\frac{a\cos i}{R_{\mathrm{s}}}\right)^2}\right)
    - a_2\left(1-\sqrt{1-\left(\frac{a\cos i}{R_{\mathrm{s}}}\right)^2}\right)^{2}}
   \right]^{1/2} \]

\item[(D01.7b)] \hfill(5 points)
\begin{center}
\begin{tabular}{cccc}
\toprule
$i$ ($^\circ$) & $b$ & $(R_{\mathrm{p}}/R_{\mathrm{s}})_{\lambda_{\mathrm{B}}}$
  & $(R_{\mathrm{p}}/R_{\mathrm{s}})_{\lambda_{\mathrm{R}}}$\\
\midrule
83.5 & 0.96 & 0.1842 & 0.1395\\
84.0 & 0.89 & 0.1547 & 0.1316\\
84.5 & 0.81 & 0.1419 & 0.1276\\
85.0 & 0.74 & 0.1341 & 0.1251\\
85.5 & 0.66 & 0.1287 & 0.1234\\
86.0 & 0.59 & 0.1248 & 0.1221\\
86.5 & 0.52 & 0.1218 & 0.1211\\
87.0 & 0.44 & 0.1196 & 0.1204\\
87.5 & 0.37 & 0.1178 & 0.1198\\
88.0 & 0.30 & 0.1165 & 0.1194\\
88.5 & 0.22 & 0.1155 & 0.1191\\
89.0 & 0.15 & 0.1149 & 0.1189\\
89.5 & 0.07 & 0.1145 & 0.1188\\
90.0 & 0.00 & 0.1143 & 0.1187\\
\bottomrule
\end{tabular}
\end{center}

\item[(D01.7c)] \hfill(7 points)\\
One can plot the full range ($i_{\min}\le i\le 90^\circ$):
\ioaafig{2025_DA_D01-s6.pdf}
However, it is enough to plot around the point of intersection (can be
deduced from the table above), to get a better resolution for the coordinates
of the intersection point.

\item[(D01.7d)] \hfill(4 points)\\
Intersection point of the graphs of
$\left(\frac{R_{\mathrm{p}}}{R_{\mathrm{s}}}\right)$ vs $i$ for
$\lambda_{\mathrm{B}}$ and $\lambda_{\mathrm{R}}$ will give the values of both
quantities.\\
Values deduced from either graph:
\begin{align*}
\left(\frac{R_{\mathrm{p}}}{R_{\mathrm{s}}}\right) &= 0.1208\\
\Longrightarrow\ R_{\mathrm{p}} &= 0.1208\times1.2\ \mathrm{R_\odot}
\end{align*}
\[ \boxed{R_{\mathrm{p}} = 0.145\ \mathrm{R_\odot}} \]
and
\[ \boxed{i = 86.75^\circ} \]

\item[(D01.8)] \hfill(2 points)
\begin{align*}
\text{Estimated mass of the planet P }
  M_{\mathrm{p}} &= 6.93\times10^{-4}\ \mathrm{M_\odot}
  = 1.389\times10^{27}\ \mathrm{kg}\\
\text{Estimated radius of the planet P }
  R_{\mathrm{p}} &= 0.145\ \mathrm{R_\odot} = 1.009\times10^{8}\ \mathrm{m}\\
\therefore\ \text{Average density of the planet} &= 320\ \mathrm{kg\,m^{-3}}
\end{align*}
The average density of the planet P is much less than that of Jupiter (a
known gaseous planet). Hence the planet P is \fbox{GASEOUS}.
\end{description}

\solhead{9.12}{TRAPPIST-1}
% source id: 2023_OBS-PLAN_3
\begin{description}
\item[(a)]~
\begin{center}
\begin{tabular}{lll}
\toprule
i. & length of the sidereal day [h] & $221\pm5$\\
ii. & orbital period [h] & $221\pm5$ (the same)\\
iii. & length of the `solar' day [h] & infinity\\
iv. & circular orbit & YES\\
v. & obliquity (axial tilt) & $0$ ($\le1^\circ$)\\
\bottomrule
\end{tabular}
\end{center}
\textit{For determining the orbital period -- 3 points, for other quantities
-- 1 point each.}

\item[(b)]~
\begin{center}
\begin{tabular}{lcccc}
\toprule
 & b & c & d & e\\
\midrule
synodic period [h] & $43.6\pm2$ & $78.9\pm2$ & $173.5\pm2$ & $433.8\pm2$\\
maximum elongation [$^\circ$] & $17.5\pm2$ & $24\pm2$ & $37\pm2$ & $49\pm2$\\
\bottomrule
\end{tabular}
\end{center}

\item[(c)]~
\begin{center}
\begin{tabular}{lcccc}
\toprule
 & b & c & d & e\\
\midrule
orbital period [h] & $36.2\pm2$ & $58.1\pm2$ & $97.2\pm2$ & $166.4\pm2$\\
semi-major axis [tau] & $0.30\pm.02$ & $0.41\pm.02$ & $0.60\pm.03$
  & $0.75\pm.04$\\
\bottomrule
\end{tabular}
\end{center}

\item[(d)]~
\begin{center}
\begin{tabular}{ll}
\toprule
Resonance & Pair of planets\\
\midrule
3:2 & e/d and f/e\\
8:5 & c/b\\
5:3 & d/c\\
8:3 & d/b\\
4:1 & e/b\\
6:1 & f/b\\
\bottomrule
\end{tabular}
\end{center}
\textit{Each resonance for 0.5 point, apart from f/e, which is worth 1
point.}
\end{description}
